\chapter{Introduction}
\label{chap:context}

% -----------------------------------------------------------------------------
\section{Context}

Semantic segmentation assigns a semantic class label to every pixel in an
image, and is fundamental to many perception systems including autonomous
driving, robotic navigation, and medical image analysis. Unlike image-level
classification, it requires fine-grained spatial understanding, especially
when distinguishing object boundaries and subtle structures that determine
whether a model can correctly separate adjacent objects and identify key
regions in a scene. Despite significant progress in recent years, modern
segmentation models still depend heavily on large-scale pixel-level
annotations. Collecting such dense labels is expensive and time-consuming:
annotating a single image from the Cityscapes urban driving dataset takes
around 1.5 hours on average~\cite{cordts2016cityscapes}, which adds up to
over 4,460 working hours for the full training set. Building effective
segmentation systems under limited annotation budgets therefore remains a
key challenge in computer vision.

To reduce this dependence on dense annotations, Hoyer et al.~\cite{hoyer2023ijcv_sde}
proposed a framework for annotation-efficient semantic segmentation, built
around three progressive stages.

\textbf{Stage (a): Standard supervised pipeline.}
This is the baseline for the whole framework, representing the simplest
possible training setup: a subset of images is manually selected from an
unlabelled pool, annotated, and used for supervised training with standard
data augmentation. No SDE components are involved at this stage, and it
serves as the reference point for measuring the improvements introduced in
later stages.

\textbf{Stage (b): SDE-enhanced semi-supervised learning pipeline.}
This stage builds on the standard pipeline by incorporating self-supervised
depth estimation (SDE) as an auxiliary signal at three points in the
training process. First, an automatic data selection strategy based on SDE
features picks the most informative samples for annotation from the
unlabelled pool, using diversity and uncertainty criteria, without requiring
any human involvement in the active learning loop. Second, the DepthMix
augmentation scheme uses depth estimates from SDE to generate
geometrically consistent mixed training samples, combined with a Mean
Teacher mechanism to make use of unlabelled data through pseudo-labels.
Third, multi-task learning trains semantic segmentation jointly with SDE as
an auxiliary task, or alternatively transfers weights from the depth branch
to initialise the segmentation branch, so that the geometric knowledge
learned by SDE can benefit the segmentation network when labelled data is
limited.

\textbf{Stage (c): Semi-supervised domain adaptation pipeline.}
This extends the Stage~(b) setup to handle cross-domain transfer from
synthetic source domains such as GTA5 and SYNTHIA to real target domains
like Cityscapes. Two new components, Cross-Domain DepthMix and Matching
Geometry Sampling, are added to use SDE for aligning the geometric
distributions between source and target domains.

This project focuses on reproducing Stages~(a) and~(b) on Cityscapes,
which provides the foundation for the architectural extensions introduced
later. Stage~(c) is outside the scope of this work, as it tackles a
different problem that requires additional datasets, domain-specific
pipelines, and alignment strategies that are not directly relevant to the
contributions here. This project instead looks at improving multi-task
feature use within the same domain, where the relationship between depth
and segmentation representations can be examined more directly.

% -----------------------------------------------------------------------------
\section{Motivation}

Most semi-supervised segmentation methods rely on pseudo-labelling and
consistency regularisation. When labels are extremely scarce, two problems
tend to arise. First, initial pseudo-labels are often unreliable, which can
lead to confirmation bias where the model keeps reinforcing its own mistakes
and gradually degrades in performance. Second, conventional mixing
augmentations like ClassMix~\cite{olsson2021classmix} tend to ignore the
underlying geometric structure of a scene, producing synthetic images that
look unrealistic and confuse the model when it tries to learn object
boundaries.

A key observation is that scene depth is closely linked to semantic
boundaries. Road surfaces, for example, are consistently flat and sit at
ground level. Self-supervised depth estimation (SDE) can extract these depth
cues from unlabelled video sequences through geometric and photometric
consistency, with no manual annotation required~\cite{godard2017unsupervised}.
Hoyer et al.~\cite{hoyer2023ijcv_sde} showed that SDE can provide a stable
geometric signal throughout the training pipeline, enabling
geometry-consistent data mixing and better feature transfer.

That said, while PAD-Net~\cite{hoyer2023ijcv_sde} already enables some
implicit cross-task interaction through its attention mechanism, where depth
and segmentation features modulate each other's self-attention outputs, this
interaction is indirect. It operates on attention-enhanced features rather
than placing any constraint on the representations themselves. More
importantly, the depth decoder receives self-supervised training signals
from all available unlabelled frames, while the segmentation decoder is only
supervised on the small labelled subset. This asymmetry means the depth
branch ends up with much richer geometric representations, yet there is no
explicit mechanism to pass that geometric knowledge into the segmentation
feature space.

Inspired by work on cross-task consistency
learning~\cite{nakano2021crosstask,zamir2020crosstask} and multi-task
learning with partial annotations~\cite{li2022mtl}, we propose a
\textbf{Cross-Task Consistency Loss} $\mathcal{L}_\text{ct}$. The idea is
fairly straightforward: features extracted from the same scene by both
decoders should reflect compatible structural information, since pixels in
the same semantic region tend to share consistent depth values, and sharp
depth changes usually coincide with semantic boundaries. Rather than letting
the two decoders learn entirely independently, we explicitly encourage their
intermediate features to stay aligned. This way, the depth decoder can act
as a continuously available, label-free teacher for the segmentation branch,
providing extra geometric supervision precisely when labelled data is most
limited.

In terms of implementation, we explore two loss formulations: a normalised
MSE computed after $\ell_2$ normalisation, and a cosine similarity loss that
directly constrains feature directions. A lightweight $1{\times}1$ projection
layer maps depth features into the segmentation feature space, and a
stop-gradient is applied on the depth side so that $\mathcal{L}_\text{ct}$
only updates the segmentation branch and the projection layer, leaving the
depth training objective untouched. A warmup period of 5{,}000 iterations is
also applied before $\mathcal{L}_\text{ct}$ is activated, giving both
decoders time to develop meaningful representations before alignment is
enforced.

We evaluate the proposed loss against three hypotheses: (H1) it improves
segmentation accuracy compared to the reproduced Stage~(b) MTL baseline;
(H2) the improvement holds consistently across different random seeds; and
(H3) a suitable $\lambda_\text{ct}$ can be identified through a systematic
loss-scale calibration procedure.

% -----------------------------------------------------------------------------
\section{Aims and Objectives}

This project has two main goals: to reproduce the SDE-enhanced segmentation
framework proposed by Hoyer et al.~\cite{hoyer2023ijcv_sde}, and to extend
the multi-task architecture with an explicit cross-task feature alignment
loss. The specific objectives are as follows.

\begin{enumerate}

\item \textbf{Reproduce Stage~(a).} Set up and verify the original codebase,
establishing a reproducible training and evaluation pipeline on
Cityscapes~\cite{cordts2016cityscapes}. This involves reproducing the fully
supervised baseline using a shared ResNet101~\cite{he2016resnet} encoder
with an ASPP~\cite{chen2018deeplabv3plus} decoder and
U-Net~\cite{ronneberger2015unet} upsampling blocks, which serves as the
reference point for all subsequent experiments.

\item \textbf{Reproduce Stage~(b).} Building on Stage~(a), reproduce four
sets of experiments in sequence: the automatic data selection for annotation
comparison, the mixing augmentation comparison, the SDE transfer and
multi-task learning comparison, and the full framework ablation study, which
measures the individual and combined performance contributions of data
selection~(S), DepthMix~(DX), and multi-task learning~(MTL).

\item \textbf{Design, implement, and evaluate a Cross-Task Consistency
Loss.} Inspired by cross-task consistency
learning~\cite{nakano2021crosstask,zamir2020crosstask} and multi-task
learning with partial annotations~\cite{li2022mtl}, propose and implement
$\mathcal{L}_\text{ct}$ within the PAD-Net~\cite{hoyer2023ijcv_sde}
multi-task architecture. Two loss formulations are investigated: normalised
MSE and cosine similarity. A loss-scale calibration procedure followed by a
$\lambda$ sweep is used to identify the best candidates, and three
independent random seeds are used to verify whether the mIoU improvements
over the Stage~(b) MTL baseline are stable and reproducible.

\item \textbf{Evaluate the full combined framework.} Integrate
$\mathcal{L}_\text{ct}$ into three configurations: MTL with DepthMix, MTL
with automatic data selection, and MTL with both DepthMix and automatic
data selection, to assess how well the proposed loss combines with the
existing framework components and measure the overall performance of the
complete system.

\end{enumerate}
% -----------------------------------------------------------------------------